\Page{Corrigés\quad /\quad Diagnostic}{Réactiver les bons réflexes}{corr-diag}
\Sub{D1\quad /\quad Longueur d'onde}
$\lambda_0=c/\nu=3{,}00\times10^8/(5{,}00\times10^{14})=6{,}00\times10^{-7}\ \mathrm{m}=600\ \mathrm{nm}$. La puissance de dix se contrôle séparément du quotient des coefficients.
\Sub{D2\quad /\quad Angle utile}
L'angle avec la normale est $90^\circ-30^\circ=60^\circ$. C'est cet angle d'incidence qu'utilise la loi de Snell-Descartes. Le rayon réfléchi fait le même angle avec la normale.
\Sub{D3\quad /\quad Grandeur conservée}
La fréquence est conservée à une interface immobile. La vitesse et la longueur d'onde changent si les indices diffèrent : $v=c/n$ et $\lambda=\lambda_0/n$.
\Sub{D4\quad /\quad Vergence}
$f'=+0{,}10\ \mathrm{m}$, donc $V=1/f'=+10\ \delta$. La conversion en mètres est indispensable pour obtenir des dioptries.
\Sub{D5\quad /\quad Distance algébrique}
$\oa=-30\ \mathrm{cm}$. L'objet est à gauche de l'origine, dans le sens opposé à celui choisi pour l'axe.
\Sub{D6\quad /\quad Image virtuelle}
Une image virtuelle est observable et photographiable : les rayons qui semblent en provenir peuvent entrer dans une seconde lentille. Ils ne convergent toutefois pas à l'emplacement virtuel ; un écran placé là ne recueille pas cette image.
\Sub{D7\quad /\quad Petit angle}
$2^\circ=2\pi/180\simeq0{,}0349\ \mathrm{rad}$. On peut donc estimer $\tan(2^\circ)\simeq0{,}0349$. Écrire $\tan(2^\circ)\simeq2$ serait une confusion d'unité.
\Sub{D8\quad /\quad Œil réduit}
La rétine joue le rôle de capteur fixe. La vergence de la lentille équivalente varie pour que l'image d'objets à différentes distances se forme sur cette rétine.
\begin{methode}{Transformer un diagnostic en travail utile}
Pour chaque réponse erronée, noter la règle qui la corrige et inventer un exemple numérique voisin. Refaire le diagnostic une semaine plus tard sans relire les solutions.
\end{methode}

\Page{Corrigés\quad /\quad E01 à E03}{Fréquence et changement d'échelle}{corr-e01}
\sol{E01}{Lire une signature spectrale}
Le spectre comporte des raies ; il est discontinu et la source est polychromatique. Avec $\nu=c/\lambda_0$ :
\begin{center}\begin{tabular}{r l r}
\toprule $\lambda_0$ (nm)&Couleur approximative&$\nu$ (Hz)\\\midrule
486,1&bleu-vert&$6{,}172\times10^{14}$\\
589,0&jaune&$5{,}093\times10^{14}$\\
656,3&rouge&$4{,}571\times10^{14}$\\\bottomrule
\end{tabular}\end{center}
Ces trois raies peuvent provenir d'un mélange ou de plusieurs transitions. Sans comparaison à une référence spectrale et sans information sur les raies absentes, l'unicité d'un élément n'est pas établie.
\sol{E02}{Une lumière dans le verre}
Dans l'air assimilé au vide, $\lambda_0=600\ \mathrm{nm}$. Dans le verre,
\[v=\frac{c}{1{,}50}=2{,}00\times10^8\ \mathrm{m\,s^{-1}},\qquad
\lambda=\frac{600}{1{,}50}=400\ \mathrm{nm}.\]
La fréquence demeure $5{,}00\times10^{14}\ \mathrm{Hz}$. On ne doit pas renommer la composante « violette » en comparant sa longueur d'onde dans le verre à une table établie dans le vide. À la sortie du verre vers l'air, sa longueur d'onde redevient $600\ \mathrm{nm}$.
\sol{E03}{Quand les rayons ne suffisent plus}
Pour $a=10^{-3}\ \mathrm{m}$, $\lambda/a=6{,}0\times10^{-4}$. Pour $a=10^{-6}\ \mathrm{m}$, le rapport vaut $0{,}60$. L'ouverture millimétrique est grande devant la longueur d'onde ; l'ouverture micrométrique ne l'est pas suffisamment pour négliger la diffraction.
Dans le premier cas, le petit angle estimé donne sur un écran à $2{,}0\ \mathrm{m}$ une demi-largeur $y\sim D\theta=1{,}2\ \mathrm{mm}$. Une ouverture grande devant $\lambda$ ne signifie donc pas que tous les effets ondulatoires sont invisibles à toute précision.
\begin{piege}{L'erreur à éviter}
Utiliser la formule de petit angle avec $\theta=0{,}60\ \mathrm{rad}$ pour annoncer une largeur précise serait incohérent. L'énoncé fournit une estimation d'échelle, pas un profil complet de diffraction.
\end{piege}

\Page{Corrigés\quad /\quad E04 à E06}{Vérifier les angles et les seuils}{corr-e04}
\sol{E04}{Le piège des $30^\circ$}
L'incidence vaut $i=60{,}0^\circ$. La loi donne
\[r=\arcsin\!\left(\frac{\sin60{,}0^\circ}{1{,}33}\right)=40{,}6^\circ.\]
L'angle entre la direction incidente et la direction transmise vaut $i-r=19{,}4^\circ$. Pour le retour inverse, le rayon arrive dans l'eau à $40{,}6^\circ$ et ressort dans l'air à $60{,}0^\circ$.
\sol{E05}{Deux faces d'un même verre}
Les seuils sont $i_{c,a}=\arcsin(1/1{,}52)=41{,}1^\circ$ et $i_{c,e}=\arcsin(1{,}33/1{,}52)=61{,}0^\circ$. À $50{,}0^\circ$, il y a réflexion totale contre l'air, mais transmission contre l'eau. Dans ce dernier cas,
\[r=\arcsin\!\left(\frac{1{,}52}{1{,}33}\sin50{,}0^\circ\right)=61{,}1^\circ.\]
L'intervalle strict $i_{c,a}<i<i_{c,e}$ assure les deux comportements recherchés avec un rayon transmis non rasant dans l'eau. On traite les égalités comme des cas limites.
\sol{E06}{Une lame à faces parallèles}
À l'entrée, $r=\arcsin(\sin45^\circ/1{,}50)=28{,}1^\circ$. À la sortie, l'incidence interne vaut aussi $r$ car les normales sont parallèles. La loi $n\sin r=\sin i'$ donne $i'=i$ : les directions sont parallèles.
Le segment parcouru dans la lame a pour longueur $\ell=e/\cos r$. L'angle entre ce segment et la direction incidente vaut $i-r$. Sa projection perpendiculaire donne
\[\result{s=\ell\sin(i-r)=e\frac{\sin(i-r)}{\cos r}=1{,}97\ \mathrm{mm}}.\]
\fig{corr-lame}{Le décalage $s$ est perpendiculaire à la direction incidente prolongée ; il ne se mesure pas le long de la face.}
\begin{methode}{Contrôles limites}
Si $i=0$, le rayon ne subit aucun décalage. Si le verre était remplacé par de l'air, $r=i$ et le décalage serait également nul. Ces deux tests confirment la structure de l'expression.
\end{methode}

\Page{Corrigés\quad /\quad E07 à E09}{Construire avant de conclure}{corr-e07}
\sol{E07}{Un miroir qui tourne}
L'image est le symétrique de l'objet, à $40{,}0\ \mathrm{cm}$ derrière le miroir. Elle est virtuelle et la distance objet-image vaut $80{,}0\ \mathrm{cm}$. Pour une rotation du miroir de $\beta$, la normale tourne aussi de $\beta$. Avec une direction incidente fixe, la direction réfléchie, symétrique de l'incidente par rapport à la normale, tourne de $2\beta$, soit $6{,}0^\circ$ en valeur absolue.
Une photographie reste possible car les rayons réfléchis atteignent l'objectif, qui en forme une image réelle sur son capteur.
\fig{miroir}{Les prolongements se rencontrent derrière le miroir ; aucun rayon réfléchi n'y pénètre.}
\sol{E08}{Assez net pour quel capteur ?}
Avec le critère fourni, $12>5$ : le premier détecteur ne juge pas le point net ; $12<20$ : le second le juge net. La tache demeure étendue dans les deux cas.
Gauss impose des rayons proches de l'axe et peu inclinés. Le système est alors approximativement stigmatique et aplanétique. Pour $3{,}0^\circ=0{,}05236\ \mathrm{rad}$, $\tan\theta=0{,}05241$ ; l'erreur relative est $9{,}15\times10^{-4}$, soit $0{,}0915\,\%$.
\sol{E09}{Reconnaître une lentille}
$L_1$ est convergente, $V_1=+12{,}5\ \delta$ ; $L_2$ est divergente, $V_2=-5{,}0\ \delta$. Dans les deux cas $F$ a pour coordonnée $-f'$ et $F'$ pour coordonnée $f'$. Le rayon par $O$ n'est pas dévié ; le rayon parallèle ressort par $F'$ ou son prolongement. Pour une direction inclinée, le foyer secondaire image est l'intersection du rayon par $O$ et du plan focal image.

\Page{Corrigés\quad /\quad E10 à E12}{Une même loi et trois natures}{corr-e10}
\sol{E10}{Une projection réduite}
$\oa=-30{,}0\ \mathrm{cm}$. On calcule $1/q=1/10-1/30=1/15$, donc $q=+15{,}0\ \mathrm{cm}$. Puis $\gamma=15/(-30)=-0{,}500$ et $\overline{A'B'}=-0{,}600\ \mathrm{cm}$. L'image est réelle, renversée et réduite. Deux rayons depuis $B$, l'un parallèle puis passant par $F'$, l'autre par $O$, se coupent à droite.
\sol{E11}{La même lentille devient une loupe}
Cette fois $1/q=1/10-1/6=-1/15$, d'où $q=-15{,}0\ \mathrm{cm}$ et $\gamma=+2{,}50$. L'image est virtuelle, droite et agrandie. Les rayons émergents divergent ; seuls leurs prolongements se coupent à gauche de la lentille.
\fig{corr-loupe}{Construction de E11 : l'objet est à $0{,}60f'$ de $O$ et l'image virtuelle à $1{,}50f'$ du même côté.}
\sol{E12}{Une divergente en autonomie}
$1/q=-1/12-1/36=-1/9$, donc $q=-9{,}00\ \mathrm{cm}$ et $\gamma=+0{,}250$. L'image est virtuelle, droite et réduite.
Plus généralement, avec une longueur objet $d>0$ et $f'=-g$, $g>0$ :
\[q=-\frac{gd}{g+d}<0,\qquad\gamma=\frac{g}{g+d}\in]0,1[.\]
Ces signes démontrent les trois propriétés sans se limiter au cas numérique.
\begin{piege}{Un contrôle de signe vaut mieux qu'une intuition isolée}
Si le calcul de E11 donnait $q>0$ ou celui de E12 donnait $\gamma<0$, il faudrait reprendre l'insertion de $\oa$ dans Descartes. Le mot « distance » ne permet pas de supprimer les signes.
\end{piege}

\Page{Corrigés\quad /\quad E13 à E15}{Passer du foyer à l'écran}{corr-e13}
\sol{E13}{Une étoile hors de l'axe}
Le faisceau incident est parallèle ; l'image est dans le plan focal image, à l'intersection du rayon passant par $O$ et de ce plan. Avec $\theta=0{,}500\pi/180=8{,}73\times10^{-3}\ \mathrm{rad}$ :
\[y'\simeq f'\theta=200\times8{,}73\times10^{-3}=+1{,}75\ \mathrm{mm}.\]
En plaçant une source au point ainsi construit et en inversant le sens de propagation, on retrouve le faisceau parallèle initial parcouru en sens inverse. Cette précision sur le sens est essentielle.
\sol{E14}{Un objet très proche du foyer}
$\overline{FA}=-1{,}00\ \mathrm{mm}$. Newton donne $\overline{F'A'}=-25/(-1)=+25{,}0\ \mathrm{mm}$, donc $\oap=30{,}0\ \mathrm{mm}$. L'objet est à $\oa=-6{,}00\ \mathrm{mm}$, d'où $\gamma=-5{,}00$. Vérification : $1/30-1/(-6)=1/5$.
Lorsque $\overline{FA}\to0^-$, l'image s'éloigne vers $+\infty$ et $|\gamma|$ augmente sans borne dans le modèle idéal. L'image sort alors de tout écran de taille finie.
\sol{E15}{Deux positions ou aucune}
$D=90{,}0\ \mathrm{cm}>4f'=80{,}0\ \mathrm{cm}$ : deux positions existent. Les solutions de $x^2-90x+1800=0$ sont $x_1=30{,}0$ et $x_2=60{,}0\ \mathrm{cm}$. Les distances images sont respectivement $60{,}0$ et $30{,}0\ \mathrm{cm}$ ; les grandissements sont $-2{,}00$ et $-0{,}500$.
\fig{bessel}{Pour deux positions, les distances objets et images s'échangent ; les grandissements sont réciproques.}
Pour $D=75{,}0\ \mathrm{cm}$, aucune position ne convient. Pour $D=80{,}0\ \mathrm{cm}$, la position unique est $x=40{,}0\ \mathrm{cm}$ et $\gamma=-1$.

\Page{Corrigés\quad /\quad E16 à E18}{Accommodation et cadrage}{corr-e16}
\sol{E16}{L'accommodation chiffrée}
L'œil emmétrope a un punctum remotum à l'infini et un punctum proximum à $25{,}0\ \mathrm{cm}$. Le modèle donne
\[V_{\min}=\frac1{0{,}0170}=58{,}8\ \delta,\qquad
V_{\max}=\frac1{0{,}0170}+\frac1{0{,}250}=62{,}8\ \delta.\]
La variation vaut $4{,}00\ \delta$. L'accommodation agit sur la vergence tandis que la rétine reste fixe ; le déplacement d'un capteur à focale constante correspond plutôt à la mise au point du modèle photographique.
\sol{E17}{Distinguer deux détails}
$\theta_{\min}=\pi/(180\times60)=2{,}91\times10^{-4}\ \mathrm{rad}$. À $100\ \mathrm{m}$, les traits sous-tendent $s/d=0{,}0200/100=2{,}00\times10^{-4}\ \mathrm{rad}$ : le critère de résolution n'est pas satisfait.
La distance maximale est $d_{\max}=s/\theta_{\min}=68{,}8\ \mathrm{m}$. Cette conclusion suppose que le contraste et l'éclairement permettent d'atteindre la limite annoncée.
\sol{E18}{Un portrait qui ne tient pas}
Avec toutes les longueurs en millimètres,
\[q=\frac{50{,}0\times2000}{2000-50{,}0}=51{,}28\ \mathrm{mm}.\]
Le tirage supplémentaire par rapport à l'infini est $q-f'=1{,}28\ \mathrm{mm}$. Le grandissement en valeur absolue vaut $q/p=0{,}02564$, donc l'image mesure $1800\times0{,}02564=46{,}2\ \mathrm{mm}$ : elle dépasse la hauteur du capteur.
Pour cadrer entièrement la personne,
\[\frac{hf'}{p-f'}\leq h_c\quad\Longleftrightarrow\quad
p\geq f'+\frac{hf'}{h_c}
=50+\frac{1800\times50}{24}=3800\ \mathrm{mm}.\]
Il faut donc au moins $3{,}80\ \mathrm{m}$ dans ce modèle, à orientation du capteur inchangée.
\begin{methode}{Le résultat raconte une expérience}
Un faible déplacement du capteur peut suffire à changer la mise au point, mais le cadrage dépend aussi de la distance de l'objet. « Faire le point » et « faire tenir dans le cadre » sont deux opérations distinctes.
\end{methode}

\Page{Corrigés\quad /\quad E19 à E21}{Les compromis d'un instrument}{corr-e19}
\sol{E19}{Exposition et mouvement}
Comme l'exposition est proportionnelle à $T/N^2$, doubler $N$ demande de multiplier $T$ par 4 : $T'=1/125\ \mathrm{s}$. Les flous valent $vT=0{,}0200\ \mathrm{mm}=20{,}0\ \mu\mathrm{m}$ puis $80{,}0\ \mu\mathrm{m}$, soit 5 et 20 pixels.
Fermer élargit la profondeur de champ géométrique, mais la durée accrue augmente ici le flou de mouvement. Une image peut être bien mise au point et néanmoins floue à cause du déplacement.
\sol{E20}{Un point encore acceptable}
$a=f'/N=6{,}25\ \mathrm{mm}$, $q_0=50{,}505\ \mathrm{mm}$ et $q=50{,}633\ \mathrm{mm}$. La coupe du cône par le capteur donne
\[b=6{,}25\left|1-\frac{50{,}505}{50{,}633}\right|
=0{,}0158\ \mathrm{mm}=15{,}8\ \mu\mathrm{m}<30{,}0\ \mu\mathrm{m}.\]
Le point est acceptable selon le critère fourni. À $N=16$, l'ouverture et le flou géométrique sont divisés par 2 : $b=7{,}89\ \mu\mathrm{m}$.
\fig{profondeur}{La netteté acceptable se décide dans le plan du capteur, même si l'image idéale se forme ailleurs.}
\sol{E21}{Injecter dans une fibre}
$\mathrm{ON}=\sqrt{1{,}48^2-1{,}46^2}=0{,}2425$ ; dans l'air, $\alpha_{\max}=14{,}0^\circ$. Pour $\alpha=10{,}0^\circ$, Snell donne $\sin\theta=\sin10^\circ/1{,}48$, donc $\theta=6{,}74^\circ$. L'incidence à la paroi est $83{,}26^\circ$, supérieure au seuil $\arcsin(1{,}46/1{,}48)\simeq80{,}57^\circ$ : le rayon est guidé. Dans l'eau, $\alpha_{\max}=\arcsin(0{,}2425/1{,}33)=10{,}5^\circ$.

\Page{Corrigés\quad /\quad E22 à E24}{Des modèles aux résultats}{corr-e22}
\sol{E22}{Une impulsion qui s'étale}
Un trajet interne d'angle $\theta$ a pour longueur $L/\cos\theta$, d'où $t(\theta)=n_1L/(c\cos\theta)$. Le trajet axial donne $t_{\min}=4{,}93\ \mu\mathrm{s}$ ; la limite de guidage impose $\cos\theta_{\max}=n_2/n_1$. Ainsi
\[\Delta t=\frac{n_1L}{c}\left(\frac{n_1}{n_2}-1\right)=67{,}6\ \mathrm{ns}.\]
Le critère fourni donne $R_{\max}=1/(2\Delta t)=7{,}40\ \mathrm{Mbit\,s^{-1}}$. Le retard de plusieurs microsecondes décrit le temps d'attente avant réception ; les dizaines de nanosecondes d'étalement décrivent la difficulté à séparer des impulsions proches.
\sol{E23}{Deux lentilles successives}
Pour $L_1$, $\overline{O_1A}=-30\ \mathrm{cm}$ donne $\overline{O_1A_1}=+15\ \mathrm{cm}$ et $\gamma_1=-1/2$. Donc $\overline{O_2A_1}=15-20=-5\ \mathrm{cm}$ : l'image intermédiaire est un objet réel pour $L_2$.
\[\frac1{\overline{O_2A_2}}=-\frac1{10}-\frac15=-\frac3{10},\qquad
\overline{O_2A_2}=-\frac{10}{3}\ \mathrm{cm}.\]
L'image finale est virtuelle pour $L_2$, à $20-10/3=16{,}7\ \mathrm{cm}$ à droite de $O_1$. Puis $\gamma_2=(-10/3)/(-5)=2/3$ et $\gamma_{\rm total}=-1/3$. Elle est renversée et réduite par rapport à l'objet initial.
\sol{E24}{Un résultat de mesure défendable}
$f'=pq/(p+q)=20{,}0\ \mathrm{cm}$, $\gamma=-q/p=-2{,}00$ et $V=+5{,}00\ \delta$. Les dérivées partielles sont $\partial f'/\partial p=q^2/(p+q)^2$ et $\partial f'/\partial q=p^2/(p+q)^2$, donc
\[u(f')=\sqrt{(4u(p)/9)^2+(u(q)/9)^2}=0{,}0458\ \mathrm{cm}.\]
On peut présenter $f'=20{,}00\ \mathrm{cm}$ avec $u(f')=0{,}05\ \mathrm{cm}$. Avec la valeur non arrondie de l'incertitude, $z=0{,}10/\sqrt{0{,}0458^2+0{,}08^2}=1{,}08$. Les résultats sont compatibles selon le repère adopté, sans que cela exclue un biais commun.

\Page{Corrigé du défi D1}{Résoudre l'inégalité de netteté}{corr-d1}
\source{Énoncé p.\,\pageref{ex-D1}. Les distances objets $p,p_0$ sont ici des longueurs positives, différentes de $\oa$.}
\Sub{1\quad /\quad Des triangles à l'inégalité}
Au plan de la lentille, le cône a un diamètre $a$ ; son sommet est à $q$. Les triangles semblables donnent $b/a=|q-q_0|/q$. Or
\[1-\frac{q_0}{q}=1-q_0\left(\frac1{f'}-\frac1p\right)
=q_0\left(\frac1p-\frac1{p_0}\right).\]
Le critère $b\leq\varepsilon$ équivaut donc à $|1/p-1/p_0|\leq k$, avec $k=\varepsilon/(a q_0)>0$.
\Sub{2\quad /\quad Deux bornes à discuter}
\[\frac1{p_0}-k\leq\frac1p\leq\frac1{p_0}+k.
\]
La borne proche est toujours $\displaystyle p_-=(1/p_0+k)^{-1}$. Si $k<1/p_0$, la borne lointaine est $\displaystyle p_+=(1/p_0-k)^{-1}$. Si $k\geq1/p_0$, toute distance suffisamment grande convient, donc $p_+=+\infty$. On conserve le domaine physique $p>f'$.
\Sub{3\quad /\quad Valeurs numériques}
En millimètres, $a=6{,}25$, $q_0=50{,}505$ et $k=9{,}504\times10^{-5}\ \mathrm{mm^{-1}}$. On trouve
\[\result{p_-=3{,}39\ \mathrm{m}},\qquad\result{p_+=9{,}53\ \mathrm{m}}.\]
Ce sont des limites de netteté acceptable pour le diamètre choisi, et non deux plans de netteté parfaite.
\Sub{4\quad /\quad Distance de réglage limite}
La borne lointaine atteint l'infini quand $k=1/p_0$, soit
\[\frac{\varepsilon}{a}\frac{p_0-f'}{f'p_0}=\frac1{p_0}
\quad\Longleftrightarrow\quad
\result{H=f'+\frac{a f'}{\varepsilon}=f'+\frac{f'^2}{N\varepsilon}}.\]
Numériquement $H=10{,}47\ \mathrm{m}$. Les effets de diffraction, d'aberration et de contraste ne sont pas prédits par cette seule géométrie ; ils empêchent de conclure qu'une fermeture arbitrairement forte améliorerait toujours la résolution.

\Page{Corrigé du défi D2}{Un budget temporel explicite}{corr-d2}
\source{Énoncé p.\,\pageref{ex-D2}. Garder les durées en secondes pendant le calcul.}
\Sub{1\quad /\quad Étalement par unité de longueur}
Le temps du trajet axial vaut $n_1L/c$. La limite la plus lente donne $n_1^2L/(cn_2)$ ; ainsi
\[\kappa=\frac{n_1}{c}\left(\frac{n_1}{n_2}-1\right)
=6{,}758\times10^{-11}\ \mathrm{s\,m^{-1}}.\]
Sur $1{,}00\ \mathrm{km}$, $\Delta t=67{,}58\ \mathrm{ns}$ et $\tau_{\rm rec}=87{,}58\ \mathrm{ns}$. À $10{,}0\ \mathrm{Mbit\,s^{-1}}$, la période binaire est $100\ \mathrm{ns}$ alors que la condition impose au moins $175{,}16\ \mathrm{ns}$ : la liaison ne satisfait pas le cahier des charges.
\Sub{2\quad /\quad Longueur maximale}
\[\tau_0+\kappa L\leq\frac1{2R}\quad\Longrightarrow\quad
\result{L_{\max}=\frac{1/(2R)-\tau_0}{\kappa}=444\ \mathrm{m}}.\]
Si $1/(2R)=\tau_0$, seule la longueur nulle satisfait le modèle. Si $1/(2R)<\tau_0$, même une liaison sans étalement est insuffisante selon le critère donné. Il ne faut pas annoncer une « longueur maximale négative » comme un résultat physique.
\Sub{3\quad /\quad Guidage et temps de transit}
L'angle d'entrée maximal vaut $14{,}03^\circ$ : le rayon à $12{,}0^\circ$ est accepté. La loi d'entrée donne $\sin\theta=\sin12^\circ/1{,}48$, soit $\theta=8{,}08^\circ$. L'incidence à la paroi est $81{,}92^\circ$, supérieure au seuil de $80{,}57^\circ$.
Avec la longueur non arrondie $443{,}919\ \mathrm{m}$,
\[t(\theta)=\frac{1{,}48\times443{,}919}{3{,}00\times10^8\cos8{,}08^\circ}
=2{,}21\ \mu\mathrm{s}.\]
\Sub{4\quad /\quad Le compromis}
Réduire l'écart $n_1-n_2$ diminue $n_1/n_2-1$ et donc la dispersion intermodale. Cela réduit aussi $\sqrt{n_1^2-n_2^2}$ : le cône d'entrée se resserre, ce qui rend l'injection plus exigeante. Une seule amélioration chiffrée ne suffit pas à juger l'ensemble d'un dispositif.

\Page{Corrigé du défi D3}{Remonter le microscope}{corr-d3}
\source{Énoncé p.\,\pageref{ex-D3}. Les grandeurs angulaires finales sont prises en valeur absolue.}
\Sub{1\quad /\quad Partir de la sortie}
Pour une image finale à l'infini, $A_1B_1$ doit être dans le plan focal objet de l'oculaire : $A_1=F_2$. Donc $q_1=\overline{O_1A_1}=f'_1+\Delta=164\ \mathrm{mm}$. Les centres sont séparés de $f'_1+\Delta+f'_2=189\ \mathrm{mm}$.
\Sub{2\quad /\quad Position de l'objet}
En notant $p_1>0$ la distance de l'objet à l'objectif, $1/p_1=1/f'_1-1/q_1$, d'où
\[p_1=\frac{f'_1(f'_1+\Delta)}{\Delta}=4{,}10\ \mathrm{mm}.\]
Ainsi $\overline{O_1A}=-4{,}10\ \mathrm{mm}$ et $\overline{F_1A}=-0{,}100\ \mathrm{mm}$ : l'objet est juste avant le foyer objet.
\Sub{3\quad /\quad Grandissement et angle}
$\gamma_1=-q_1/p_1=-\Delta/f'_1=-40{,}0$. Pour $h=10{,}0\ \mu\mathrm{m}$, la hauteur intermédiaire vaut $-0{,}400\ \mathrm{mm}$. L'oculaire donne $\alpha'\simeq|h_1|/f'_2=0{,}0160\ \mathrm{rad}$.
\[\result{P_i=\frac{\alpha'}h=\frac{\Delta}{f'_1f'_2}=1{,}60\times10^3\ \mathrm{m^{-1}}}.\]
\fig{microscope}{L'image intermédiaire est réelle et renversée ; elle sert d'objet à l'oculaire. Le dessin conserve les conjugaisons, sans reprendre les proportions numériques du défi.}
\Sub{4\quad /\quad Interpréter le grossissement}
À l'œil nu à $d_m=0{,}250\ \mathrm{m}$, $\alpha\simeq h/d_m=4{,}00\times10^{-5}\ \mathrm{rad}$. Le grossissement en valeur absolue est $\alpha'/\alpha=P_i d_m=400$. L'œil emmétrope peut observer sans accommoder. Ce grossissement ne garantit pas la résolution : diffraction, ouverture, aberrations et contraste restent à considérer.

\Page{Corrigé du DM\quad /\quad Sur 20}{La focale et son incertitude}{corr-dm}
\source{Énoncé p.\,\pageref{ex-DM}. Barème proposé : il valorise la méthode même si une erreur numérique isolée est reportée de manière cohérente.}
\Sub{1 et 2\quad /\quad Construire la méthode\quad /\quad 7 points}
$\oa=-x$, $\oap=D-x$, donc $1/x+1/(D-x)=1/f'$. On obtient $x^2-Dx+Df'=0$ : \textbf{2 points}. Le discriminant $D(D-4f')$ impose $D\geq4f'$ : \textbf{2 points}.
Les racines sont $(D\mp\sqrt{D^2-4Df'})/2$ : \textbf{1 point}. Leur différence donne $d^2=D^2-4Df'$, puis $f'=(D^2-d^2)/(4D)$ : \textbf{2 points}.
\Sub{3 et 4\quad /\quad Exploiter les positions\quad /\quad 4 points}
$f'=160{,}0\ \mathrm{mm}$ et $x_1=(D-d)/2=200{,}0\ \mathrm{mm}$, $x_2=(D+d)/2=800{,}0\ \mathrm{mm}$ : \textbf{2 points}.
$\gamma_1=-(D-x_1)/x_1=-4{,}00$ et $\gamma_2=-0{,}250$ ; leur produit vaut 1 car les distances objets et images s'échangent : \textbf{2 points}.
\Sub{5\quad /\quad Propager l'incertitude\quad /\quad 3 points}
\[\frac{\partial f'}{\partial D}=\frac14+\frac{d^2}{4D^2}=0{,}340,
\qquad \frac{\partial f'}{\partial d}=-\frac d{2D}=-0{,}300.\]
\[u(f')=\sqrt{(0{,}340\times1{,}0)^2+(0{,}300\times2{,}0)^2}
=0{,}690\ \mathrm{mm}.\]
Une présentation adaptée est $f'=160{,}0\ \mathrm{mm}$ avec $u(f')=0{,}7\ \mathrm{mm}$. Dérivées, composition et présentation : \textbf{1 point chacune}.
\Sub{6\quad /\quad Comparer\quad /\quad 2 points}
$z=|160{,}0-158{,}5|/\sqrt{0{,}690^2+0{,}8^2}=1{,}42$. Les valeurs sont compatibles au repère $z\leq2$ ; le résultat ne prouve pas l'absence de biais commun. \textbf{1 point pour le calcul, 1 pour l'interprétation.}
\Sub{7 et 8\quad /\quad Expérimenter et critiquer\quad /\quad 4 points}
Déplacer la lentille de part et d'autre du meilleur réglage, relever l'intervalle jugé net, puis répéter indépendamment la mise au point. Une estimation d'incertitude doit préciser le modèle de distribution : \textbf{2 points}.
Exemples de limites : épaisseur non négligeable de la lentille, mauvaise référence au centre optique, défaut d'alignement, aberrations, parallaxe de lecture : \textbf{1 point par limite correctement expliquée, maximum 2}.

\Page{Corrigé du DS\quad /\quad Partie A}{Un seuil physique pour détecter l'eau}{corr-ds1}
\source{Énoncé p.\,\pageref{ex-DS}. Partie A : 8 points sur les 20 du devoir.}
\Sub{A1\quad /\quad Les lois\quad /\quad 2 points}
Les rayons sont dans le plan d'incidence. L'angle réfléchi égale l'angle incident, tous deux mesurés depuis la normale. La réfraction vérifie $n_v\sin i=n_{\rm ext}\sin r$. Attribution : \textbf{1 point} pour réflexion et coplanarité, \textbf{1 point} pour réfraction avec identification des angles et milieux.
\Sub{A2\quad /\quad Les seuils\quad /\quad 2 points}
L'existence de $r$ exige $(n_v/n_{\rm ext})\sin i\leq1$. Comme $n_v>n_{\rm ext}$, la réflexion totale commence strictement au-delà de $i_c=\arcsin(n_{\rm ext}/n_v)$. On obtient
\[\result{i_{c,a}=41{,}8^\circ},\qquad\result{i_{c,e}=62{,}5^\circ}.\]
\textbf{1 point} pour la démonstration et \textbf{1 point} pour les deux valeurs.
\Sub{A3\quad /\quad À l'incidence imposée\quad /\quad 2 points}
À $50{,}0^\circ$, le verre-air est en réflexion totale. Le verre-eau admet un rayon transmis :
\[r=\arcsin\!\left(\frac{1{,}50}{1{,}33}\sin50{,}0^\circ\right)=59{,}8^\circ.\]
\textbf{1 point} pour les comportements et \textbf{1 point} pour $r$ avec son expression.
\fig{capteur-pluie}{Lorsqu'un rayon transmis apparaît, une partie de l'énergie peut quitter le verre.}
\Sub{A4\quad /\quad Une information exploitable\quad /\quad 2 points}
Le domaine $41{,}8^\circ<i<62{,}5^\circ$, en gardant les seuils non arrondis pour la décision, donne réflexion totale à sec et transmission non rasante dans l'eau : \textbf{1 point}. Dans le modèle idéal sans absorption, l'énergie transmise réduit l'énergie réfléchie reçue par le détecteur : \textbf{1 point}.
\begin{piege}{Ce que l'on ne peut pas chiffrer ici}
La loi de Snell-Descartes ne donne pas le rapport des puissances réfléchies. Le devoir demande une explication qualitative, pas un pourcentage de signal inventé.
\end{piege}

\Page{Corrigé du DS\quad /\quad Parties B et C}{Projection et transmission}{corr-ds2}
\Sub{B1 et B2\quad /\quad L'image\quad /\quad 4 points}
$\oa=-300\ \mathrm{mm}$ et $1/\oap-1/\oa=1/f'$ : \textbf{1 point}. Donc $\oap=+150\ \mathrm{mm}$ : \textbf{1 point} ; $\gamma=-0{,}500$ et $h'=-6{,}00\ \mathrm{mm}$ : \textbf{1 point} ; image réelle, renversée et réduite : \textbf{1 point}.
\Sub{B3 et B4\quad /\quad Deux constructions\quad /\quad 4 points}
$450>4\times100$ : deux positions sont possibles. Les solutions sont $x=(450\mp150)/2$, soit $150$ et $300\ \mathrm{mm}$ depuis l'objet : \textbf{2 points} pour condition et positions.
Construction initiale : rayon de $B$ parallèle à l'axe puis passant par $F'$ ; rayon de $B$ passant par $O$ sans déviation. Ils se coupent à $150\ \mathrm{mm}$ à droite : \textbf{1 point}. Dans l'autre position, $q=300\ \mathrm{mm}$ et $p=150\ \mathrm{mm}$, donc $\gamma=-2{,}00$ et $h'=-24{,}0\ \mathrm{mm}$ : \textbf{1 point}.
\Sub{C1\quad /\quad Guidage\quad /\quad 2 points}
L'incidence à la paroi est $i=\pi/2-\theta$. La réflexion totale impose $\cos\theta>n_2/n_1$. À l'entrée $\sin\alpha=n_1\sin\theta$, d'où $\sin^2\alpha<n_1^2-n_2^2$ : \textbf{1 point}. La limite donne $\alpha_{\max}=\arcsin\sqrt{1{,}48^2-1{,}46^2}=14{,}0^\circ$ : \textbf{1 point}.
\Sub{C2\quad /\quad Étalement et débit\quad /\quad 2 points}
Le temps interne est $n_1L/(c\cos\theta)$ ; les extrêmes sont $n_1L/c$ et $n_1^2L/(cn_2)$. Leur différence donne
\[\Delta t=\frac{n_1L}{c}\left(\frac{n_1}{n_2}-1\right)=33{,}8\ \mathrm{ns}.\]
Démonstration et valeur : \textbf{1 point}. Le critère $T\geq2\Delta t$ donne $R\leq1/(2\Delta t)=14{,}8\ \mathrm{Mbit\,s^{-1}}$ : \textbf{1 point}.
\begin{methode}{Analyser sa note par capacité}
Une note totale masque parfois une fragilité locale. Reprendre séparément : angles et seuils (A), signes et construction (B), chaîne de raisonnement (C). Un schéma inexpliqué ou un résultat sans unité ne constitue pas une démonstration.
\end{methode}
